% GENERATED FILE. Edit canonical lessons, then run npm run generate:books.
% canonical-source-hash: fnv1a64:2e2a128069e9fc26
% canonical-lessons: SA-C460-svabhavah, SA-C460-sabhyah, SA-C460-pravartate, SA-C460-gunah, SA-C460-dustah

\chapter{Character and Conduct}
\label{ch:sa-character-and-conduct}
\hlchaptermodality{460}

\begin{chapteropening}
\textbf{By the end of this chapter:} \emph{I can describe someone's character.}

Complete the last lesson of chapter 460: I can describe someone's character.
\end{chapteropening}

\section[\texorpdfstring{\sk{स्वभावः} (svabhāvaḥ)}{स्वभावः (svabhāvaḥ)}]{\sk{स्वभावः} (svabhāvaḥ) — temperament}

\label{lesson:SA-C460-svabhavah}



\begin{quote}
Before the new one: say the Sanskrit for regard, then the Sanskrit for self-respect.
\end{quote}

\subsection*{You'll want to know: \sk{स्वभावः}}
\textbf{\sk{स्वभावः}} — \emph{svabhāvaḥ} — \textquotedblleft{}temperament\textquotedblright{}.

\begin{quote}
\sk{मीरायाः} \sk{स्वभावः} \sk{शान्तः} \sk{अस्ति।}
\end{quote}

\begin{quote}
\emph{mīrāyāḥ svabhāvaḥ śāntaḥ asti.}
\end{quote}

\begin{quote}
Mira has a calm temperament.
\end{quote}

A masculine noun.

\begin{cousinweb}
\textbf{\sk{स्वभावः}} (\emph{svabhāvaḥ}) is built from \emph{sva}, \textquotedblleft{}own\textquotedblright{}, and \emph{bhāva}, \textquotedblleft{}being\textquotedblright{}, from the root √bhū: one's own way of being; MW gives \textquotedblleft{}natural state, innate disposition, nature\textquotedblright{} (MW).
\end{cousinweb}

\begin{grammarlens}[title={What you've built}]
One more word for describing what you lived and felt.
\end{grammarlens}

\subsection*{Your turn}
\begin{itemize}
\raggedright
  \item \emph{Say it:} \emph{svabhāvaḥ}
  \item \emph{Say it:} \emph{svabhāvaḥ}, once more
  \item \emph{Say it:} \emph{mīrāyāḥ svabhāvaḥ śāntaḥ asti.} — Mira has a calm temperament.
  \item \emph{Recall:} say the Sanskrit for regard, then the Sanskrit for self-respect, then say \emph{svabhāvaḥ} again
\end{itemize}

\begin{tcolorbox}[breakable,colback=teal!4,colframe=teal!35!black,title={Before you move on}]
What does \sk{स्वभावः} mean? (\textquotedblleft{}Temperament\textquotedblright{}.) Say it once more.
\end{tcolorbox}
\section[\texorpdfstring{\sk{सभ्यः} (sabhyaḥ)}{सभ्यः (sabhyaḥ)}]{\sk{सभ्यः} (sabhyaḥ) — refined; courteous}

\label{lesson:SA-C460-sabhyah}



\begin{quote}
Before the new one: say the Sanskrit for self-respect, then the Sanskrit for temperament.
\end{quote}

\subsection*{You'll want to know: \sk{सभ्यः}}
\textbf{\sk{सभ्यः}} — \emph{sabhyaḥ} — \textquotedblleft{}refined; courteous\textquotedblright{}.

\begin{quote}
\sk{सभ्यः} \sk{जनः} \sk{कटु} \sk{न} \sk{वदति।}
\end{quote}

\begin{quote}
\emph{sabhyaḥ janaḥ kaṭu na vadati.}
\end{quote}

\begin{quote}
A refined person does not speak harshly.
\end{quote}

\begin{cousinweb}
\textbf{\sk{सभ्यः}} (\emph{sabhyaḥ}) is built on \textbf{\sk{सभा}}, \textquotedblleft{}an assembly, a hall\textquotedblright{}: fit for an assembly; MW gives \textquotedblleft{}suitable to good society, courteous, polite, refined, civilized\textquotedblright{} (MW).
\end{cousinweb}

\begin{grammarlens}[title={What you've built}]
One more word for describing what you lived and felt.
\end{grammarlens}

\subsection*{Your turn}
\begin{itemize}
\raggedright
  \item \emph{Say it:} \emph{sabhyaḥ}
  \item \emph{Say it:} \emph{sabhyaḥ}, once more
  \item \emph{Say it:} \emph{sabhyaḥ janaḥ kaṭu na vadati.} — A refined person does not speak harshly.
  \item \emph{Recall:} say the Sanskrit for self-respect, then the Sanskrit for temperament, then say \emph{sabhyaḥ} again
\end{itemize}

\begin{tcolorbox}[breakable,colback=teal!4,colframe=teal!35!black,title={Before you move on}]
What does \sk{सभ्यः} mean? (\textquotedblleft{}Refined; courteous\textquotedblright{}.) Say it once more.
\end{tcolorbox}
\section[\texorpdfstring{\sk{प्रवर्तते} (pravartate)}{प्रवर्तते (pravartate)}]{\sk{प्रवर्तते} (pravartate) — conducts himself; proceeds}

\label{lesson:SA-C460-pravartate}



\begin{quote}
Before the new one: say the Sanskrit for temperament, then the Sanskrit for refined.
\end{quote}

\subsection*{You'll want to know: \sk{प्रवर्तते}}
\textbf{\sk{प्रवर्तते}} — \emph{pravartate} — \textquotedblleft{}conducts himself; proceeds\textquotedblright{}.

\begin{quote}
\sk{सः} \sk{विनयेन} \sk{प्रवर्तते।}
\end{quote}

\begin{quote}
\emph{saḥ vinayena pravartate.}
\end{quote}

\begin{quote}
He conducts himself with modesty.
\end{quote}

Past (he …): \sk{सः} \sk{विनयेन} \sk{प्रवृत्तवान्}, \emph{saḥ vinayena pravṛttavān}.

\begin{cousinweb}
\textbf{\sk{प्रवर्तते}} (\emph{pravartate}) is built from \emph{pra}, \textquotedblleft{}forth\textquotedblright{}, and the root √vṛt, \textquotedblleft{}to turn\textquotedblright{}, the root of \textbf{\sk{वर्तते}}; MW gives \textquotedblleft{}to proceed, act\textquotedblright{}, and its ppp \emph{pravṛtta} \textquotedblleft{}engaged in, occupied with\textquotedblright{} (MW).
\end{cousinweb}

\begin{grammarlens}[title={What you've built}]
One more word for describing what you lived and felt.
\end{grammarlens}

\subsection*{Your turn}
\begin{itemize}
\raggedright
  \item \emph{Say it:} \emph{pravartate}
  \item \emph{Say it:} \emph{pravartate}, once more
  \item \emph{Say it:} \emph{saḥ vinayena pravartate.} — He conducts himself with modesty.
  \item \emph{Recall:} say the Sanskrit for temperament, then the Sanskrit for refined, then say \emph{pravartate} again
\end{itemize}

\begin{tcolorbox}[breakable,colback=teal!4,colframe=teal!35!black,title={Before you move on}]
What does \sk{प्रवर्तते} mean? (\textquotedblleft{}Conducts himself; proceeds\textquotedblright{}.) Say it once more.
\end{tcolorbox}
\section[\texorpdfstring{\sk{गुणः} (guṇaḥ)}{गुणः (guṇaḥ)}]{\sk{गुणः} (guṇaḥ) — a virtue}

\label{lesson:SA-C460-gunah}



\begin{quote}
Before the new one: say the Sanskrit for refined, then the Sanskrit for conducts himself.
\end{quote}

\subsection*{You'll want to know: \sk{गुणः}}
\textbf{\sk{गुणः}} — \emph{guṇaḥ} — \textquotedblleft{}a virtue\textquotedblright{}.

\begin{quote}
\sk{दया} \sk{महान्} \sk{गुणः} \sk{अस्ति।}
\end{quote}

\begin{quote}
\emph{dayā mahān guṇaḥ asti.}
\end{quote}

\begin{quote}
Compassion is a great virtue.
\end{quote}

A masculine noun.

\begin{cousinweb}
\textbf{\sk{गुणः}} (\emph{guṇaḥ}) is traced to the root √grah, \textquotedblleft{}to seize\textquotedblright{}, in the Uṇādi list that MW follows; its first sense is \textquotedblleft{}a strand of a cord\textquotedblright{}, and \textquotedblleft{}a good quality, virtue\textquotedblright{} comes later (MW).
\end{cousinweb}

\begin{grammarlens}[title={What you've built}]
One more word for describing what you lived and felt.
\end{grammarlens}

\subsection*{Your turn}
\begin{itemize}
\raggedright
  \item \emph{Say it:} \emph{guṇaḥ}
  \item \emph{Say it:} \emph{guṇaḥ}, once more
  \item \emph{Say it:} \emph{dayā mahān guṇaḥ asti.} — Compassion is a great virtue.
  \item \emph{Recall:} say the Sanskrit for refined, then the Sanskrit for conducts himself, then say \emph{guṇaḥ} again
\end{itemize}

\begin{tcolorbox}[breakable,colback=teal!4,colframe=teal!35!black,title={Before you move on}]
What does \sk{गुणः} mean? (\textquotedblleft{}A virtue\textquotedblright{}.) Say it once more.
\end{tcolorbox}
\section[\texorpdfstring{\sk{दुष्टः} (duṣṭaḥ)}{दुष्टः (duṣṭaḥ)}]{\sk{दुष्टः} (duṣṭaḥ) — wicked; corrupt}

\label{lesson:SA-C460-dustah}



\begin{quote}
Before the new one: say the Sanskrit for conducts himself, then the Sanskrit for a virtue.
\end{quote}

\subsection*{You'll want to know: \sk{दुष्टः}}
\textbf{\sk{दुष्टः}} — \emph{duṣṭaḥ} — \textquotedblleft{}wicked; corrupt\textquotedblright{}.

\begin{quote}
\sk{दुष्टः} \sk{बालः} \sk{शुनकं} \sk{ताडयति।}
\end{quote}

\begin{quote}
\emph{duṣṭaḥ bālaḥ śunakaṃ tāḍayati.}
\end{quote}

\begin{quote}
The wicked boy beats the dog.
\end{quote}

\begin{cousinweb}
\textbf{\sk{दुष्टः}} (\emph{duṣṭaḥ}) is built on the root √duṣ, \textquotedblleft{}to become bad, to spoil\textquotedblright{}: spoilt; MW gives \textquotedblleft{}spoilt, corrupted\textquotedblright{} and \textquotedblleft{}bad, wicked\textquotedblright{} (MW).
\end{cousinweb}

\begin{grammarlens}[title={What you've built}]
One more word for describing what you lived and felt.
\end{grammarlens}

\subsection*{Your turn}
\begin{itemize}
\raggedright
  \item \emph{Say it:} \emph{duṣṭaḥ}
  \item \emph{Say it:} \emph{duṣṭaḥ}, once more
  \item \emph{Say it:} \emph{duṣṭaḥ bālaḥ śunakaṃ tāḍayati.} — The wicked boy beats the dog.
  \item \emph{Recall:} say the Sanskrit for conducts himself, then the Sanskrit for a virtue, then say \emph{duṣṭaḥ} again
\end{itemize}

\begin{tcolorbox}[breakable,colback=teal!4,colframe=teal!35!black,title={Before you move on}]
What does \sk{दुष्टः} mean? (\textquotedblleft{}Wicked; corrupt\textquotedblright{}.) Say it once more.
\end{tcolorbox}
